\section{Form: Tipe input, select, textarea, dan validasi dasar}

HTML5 memperkenalkan banyak tipe input baru yang melampaui \code{type="text"} dan \code{type="password"} tradisional. Tipe-tipe ini memberikan petunjuk kepada browser tentang format data yang diharapkan, sehingga browser dapat menampilkan kontrol yang sesuai---misalnya kalender untuk \code{type="date"}, slider untuk \code{type="range"}, atau keyboard numerik di perangkat mobile untuk \code{type="tel"}. Penggunaan tipe yang tepat meningkatkan pengalaman pengguna dan mengaktifkan validasi bawaan browser tanpa memerlukan JavaScript tambahan \cite{mdnHtmlIntro}.

\begin{table}[htbp]
\centering
\begin{tabular}{|P{2.2cm}|P{3.5cm}|P{3cm}|P{3cm}|}
\hline
\textbf{Tipe Input} & \textbf{Fungsi} & \textbf{Contoh Nilai} & \textbf{Kontrol UI} \\
\hline
\code{text} & Teks satu baris & ``Budi Santoso'' & Text field \\
\hline
\code{email} & Alamat email & ``budi@mail.com'' & Validasi format @ \\
\hline
\code{password} & Kata sandi & ``********'' & Teks tersembunyi \\
\hline
\code{number} & Angka & ``25'' & Spinner +/-- \\
\hline
\code{date} & Tanggal & ``2026-01-15'' & Date picker \\
\hline
\code{time} & Waktu & ``14:30'' & Time picker \\
\hline
\code{url} & URL & ``https://...'' & Validasi format URL \\
\hline
\code{tel} & Nomor telepon & ``08123456789'' & Keyboard telepon \\
\hline
\code{range} & Slider rentang & ``50'' & Slider \\
\hline
\code{color} & Warna & ``\#ff0000'' & Color picker \\
\hline
\code{checkbox} & Pilihan ganda & on/off & Kotak centang \\
\hline
\code{radio} & Pilihan tunggal & satu dari grup & Tombol radio \\
\hline
\code{file} & Unggah file & (path file) & Tombol browse \\
\hline
\code{hidden} & Data tersembunyi & (tidak terlihat) & Tidak ada UI \\
\hline
\end{tabular}
\caption{Tipe-tipe input HTML5 dan karakteristiknya}
\end{table}

Elemen \code{<select>} digunakan untuk membuat daftar pilihan dropdown. Setiap pilihan didefinisikan dengan \code{<option>} dan dapat dikelompokkan menggunakan \code{<optgroup>}. Atribut \code{multiple} memungkinkan pengguna memilih lebih dari satu opsi (biasanya dengan Ctrl+klik). Untuk teks panjang seperti komentar atau deskripsi, gunakan \code{<textarea>} yang menyediakan area input multi-baris. Ukuran awal textarea dapat diatur dengan atribut \code{rows} dan \code{cols}, dan pengguna umumnya dapat mengubah ukurannya secara interaktif \cite{webdevHtml}.

Validasi form di sisi klien (client-side) membantu pengguna mengisi data dengan benar sebelum dikirim ke server. HTML5 menyediakan beberapa atribut validasi bawaan: \code{required} mewajibkan field diisi, \code{minlength} dan \code{maxlength} membatasi panjang karakter, \code{min} dan \code{max} membatasi rentang numerik, dan \code{pattern} mencocokkan input dengan pola regular expression. Browser akan menampilkan pesan error otomatis saat validasi gagal. Perlu diingat bahwa validasi klien hanya untuk kenyamanan pengguna; \textbf{validasi server tetap wajib} untuk keamanan karena data dari browser dapat dimanipulasi \cite{whatwgHtml}.

\begin{konsep}
\textbf{Validasi Client-side vs Server-side.}
\begin{itemize}
  \item \textbf{Client-side} (HTML/JS): Memberikan umpan balik instan kepada pengguna; mengurangi request ke server yang tidak perlu; tetapi dapat dilewati (bypass) oleh pengguna.
  \item \textbf{Server-side} (PHP/Node.js/dll.): Menjamin keamanan data; tidak dapat dilewati; wajib untuk semua form yang menerima data publik.
\end{itemize}
Pendekatan terbaik: gunakan keduanya secara bersamaan. Client-side untuk UX, server-side untuk keamanan \cite{mdnHtmlIntro}.
\end{konsep}

Elemen \code{<fieldset>} dan \code{<legend>} digunakan untuk mengelompokkan kontrol form yang saling terkait. \code{<fieldset>} membungkus sekumpulan field dan secara visual menampilkan bingkai, sedangkan \code{<legend>} memberikan judul untuk grup tersebut. Pengelompokan ini sangat penting untuk aksesibilitas---khususnya pada grup radio button dan checkbox yang membutuhkan konteks agar pembaca layar dapat menjelaskan pilihan yang tersedia \cite{webdevAccessibility}.

\begin{webcode}[language=HTML, caption={Form pendaftaran dengan fieldset, select, dan validasi}]
<form action="/daftar" method="post">
  <fieldset>
    <legend>Data Pribadi</legend>
    <label for="nama">Nama:</label>
    <input type="text" id="nama" name="nama"
           required minlength="3" />
    <label for="email">Email:</label>
    <input type="email" id="email" name="email" required />
    <label for="tgl">Tanggal Lahir:</label>
    <input type="date" id="tgl" name="tgl_lahir" />
  </fieldset>

  <fieldset>
    <legend>Informasi Akademik</legend>
    <label for="jurusan">Jurusan:</label>
    <select id="jurusan" name="jurusan" required>
      <option value="">-- Pilih Jurusan --</option>
      <optgroup label="Fakultas Teknik">
        <option value="if">Teknik Informatika</option>
        <option value="si">Sistem Informasi</option>
      </optgroup>
      <optgroup label="Fakultas Ekonomi">
        <option value="mn">Manajemen</option>
        <option value="ak">Akuntansi</option>
      </optgroup>
    </select>
    <label for="bio">Motivasi:</label>
    <textarea id="bio" name="bio" rows="4"
              maxlength="500"
              placeholder="Ceritakan motivasi Anda...">
    </textarea>
  </fieldset>

  <fieldset>
    <legend>Jenis Kelamin</legend>
    <label>
      <input type="radio" name="gender" value="L" required />
      Laki-laki
    </label>
    <label>
      <input type="radio" name="gender" value="P" />
      Perempuan
    </label>
  </fieldset>

  <label>
    <input type="checkbox" name="setuju" required />
    Saya menyetujui syarat dan ketentuan
  </label>

  <button type="submit">Daftar</button>
</form>
\end{webcode}

